% ============================================================
%  Rapport d'etat — Pipeline de transcription broadcast -> SRT
%  HACA / Direction des Systemes d'Information
%  Auteur : Marwane ElBaraka
%  Compiler : pdflatex (2 passes) ou Overleaf
% ============================================================
\documentclass[11pt, a4paper]{article}

% ---------- Encodage & langue ----------
\usepackage[utf8]{inputenc}
\usepackage[T1]{fontenc}
\usepackage[french]{babel}

% ---------- Mise en page ----------
\usepackage[top=2cm, bottom=2cm, left=2.4cm, right=2.4cm]{geometry}
\usepackage{parskip}
\setlength{\parindent}{0pt}
\usepackage{lmodern}
\usepackage{microtype}

% ---------- Couleurs ----------
\usepackage[dvipsnames, table]{xcolor}
\definecolor{hacablue}{RGB}{0,70,127}
\definecolor{goodgreen}{RGB}{0,150,80}
\definecolor{warnorange}{RGB}{200,110,0}
\definecolor{darkgray}{RGB}{80,80,80}

% ---------- Tableaux ----------
\usepackage{booktabs}
\usepackage{array}
\usepackage{tabularx}
\usepackage{colortbl}

% ---------- Listes ----------
\usepackage{enumitem}
\setlist[itemize]{noitemsep, topsep=2pt, leftmargin=1.2em}
\setlist[enumerate]{noitemsep, topsep=2pt, leftmargin=1.4em}

% ---------- Code ----------
\usepackage{listings}
\lstset{basicstyle=\ttfamily\footnotesize, backgroundcolor=\color{gray!8},
  frame=single, framerule=0.3pt, breaklines=true, columns=flexible, xleftmargin=2pt}

% ---------- Titres ----------
\usepackage{titlesec}
\titleformat{\section}{\large\bfseries\color{hacablue}}{\thesection.}{0.5em}{}[\vspace{-3pt}\rule{\linewidth}{0.4pt}]
\titleformat{\subsection}{\normalsize\bfseries\color{hacablue!80!black}}{\thesubsection.}{0.4em}{}
\titlespacing*{\section}{0pt}{8pt}{4pt}
\titlespacing*{\subsection}{0pt}{6pt}{2pt}

% ---------- Boites ----------
\usepackage{tcolorbox}
\tcbuselibrary{skins, breakable}
\newtcolorbox{resume}{colback=hacablue!5, colframe=hacablue!60,
  fonttitle=\bfseries, title=Resume executif, boxrule=0.5pt, arc=2pt}

% ---------- En-tete / pied ----------
\usepackage{fancyhdr}
\pagestyle{fancy}
\fancyhf{}
\fancyhead[L]{\small\textcolor{darkgray}{Rapport d'etat — Transcription broadcast}}
\fancyhead[R]{\small\textcolor{darkgray}{HACA / DSI}}
\fancyfoot[C]{\small\thepage}
\renewcommand{\headrulewidth}{0.4pt}

\usepackage[hidelinks, colorlinks=true, linkcolor=hacablue, urlcolor=hacablue]{hyperref}

% ============================================================
\begin{document}

\begin{center}
  {\LARGE\bfseries\color{hacablue} Pipeline de transcription broadcast $\rightarrow$ SRT}\\[3pt]
  {\large pour contenus en darija, arabe standard et francais}\\[8pt]
  {\normalsize \textbf{Rapport d'etat}}\\[4pt]
  {\small Stage d'initiation — HACA / Direction des Systemes d'Information}\\[1pt]
  {\small Stagiaire : Marwane ElBaraka \quad|\quad Encadrante : Mme.~Amina Moujtahid \quad|\quad Date : \today}
\end{center}

\vspace{2pt}\rule{\linewidth}{0.8pt}\vspace{6pt}

\begin{resume}
Cet outil convertit des fichiers \textbf{audio/video} de diffusion en sous-titres
\textbf{\texttt{.srt}} standard, pour des contenus en \textbf{darija marocain, arabe
standard moderne (MSA) et francais}, y compris les fichiers qui \textbf{alternent} entre
ces langues. Il est bati sur \textbf{faster-whisper} (Whisper \texttt{large-v3}) avec une
\textbf{detection de langue par segment}, ce qui permet de transcrire nativement une
emission multilingue. Deux pipelines coexistent : un pipeline leger (faster-whisper) et
un pipeline \textbf{WhisperX} qui ajoute l'alignement mot-a-mot et la diarisation des
locuteurs. Un \textbf{adaptateur LoRA Darija} ameliore la qualite sur l'arabe marocain.
La sortie \texttt{.srt} est directement consommable par le benchmark de tonalite HACA :
les deux projets ne partagent que le format de fichier.
\end{resume}

% ============================================================
\section{Objectif et positionnement}
% ============================================================

Cet outil constitue la \emph{moitie amont} de la chaine d'analyse media : le benchmark de
tonalite consommait deja des \texttt{.srt} mais il manquait un moyen d'en \emph{produire}
a partir de diffusions brutes. Les deux projets sont volontairement separes ; ils ne
partagent rien d'autre que le format \texttt{.srt}.

Le defi principal est le \textbf{code-switching} : une emission marocaine peut s'ouvrir
en francais, basculer en darija et citer du MSA. Whisper standard ne detecte qu'\emph{une}
langue sur les $\sim$30 premieres secondes et l'applique a tout le fichier — ce qui est
faux pour la majorite du contenu. La detection par segment resout ce probleme.

% ============================================================
\section{Architecture du pipeline}
% ============================================================

\begin{lstlisting}
audio -> decode 16kHz mono -> VAD (Silero) -> chunks ~25s (coupe au silence)
      -> detection langue par chunk (allow-list ar,fr,en)
      -> transcription par chunk (condition_on_previous_text=False)
      -> offset des timestamps -> ecriture .srt
\end{lstlisting}

\begin{enumerate}
  \item \textbf{Decodage} en 16 kHz mono (PyAV, gere aussi la video).
  \item \textbf{VAD Silero} : detection des zones de parole, le silence est ecarte.
  \item \textbf{Regroupement} en segments d'au plus 25 s, coupes \emph{uniquement au silence} (jamais au milieu d'un mot).
  \item \textbf{Detection de langue par chunk}, contrainte a une \textbf{allow-list} (\texttt{ar,fr,en}) ; hors-liste $\rightarrow$ repli sur \texttt{ar} (evite que du darija bruite soit pris pour du persan/ourdou).
  \item \textbf{Transcription} de chaque chunk dans sa langue, avec \texttt{condition\_on\_previous\_text=False} pour eviter qu'une hallucination ne se propage.
  \item \textbf{Decalage des timestamps} selon le debut absolu du chunk, puis ecriture du \texttt{.srt} (index, \texttt{HH:MM:SS,mmm}, UTF-8).
\end{enumerate}

% ============================================================
\section{Deux pipelines}
% ============================================================

\begin{center}
\small
\begin{tabularx}{\linewidth}{lXX}
\toprule
 & \textbf{A — \texttt{transcribe.py}} & \textbf{B — \texttt{transcribe\_whisperx.py}} \\
\midrule
Moteur & faster-whisper (CTranslate2) & WhisperX (memes chunks) \\
Dependances & legeres & lourdes (pyannote, ONNX) \\
Atouts & transcription simple, batch, tests CPU, production T4 & alignement mot-a-mot (wav2vec2) + diarisation locuteurs \\
Quand l'utiliser & besoin uniquement de \texttt{.srt} & besoin de locuteurs ou timings precis \\
\bottomrule
\end{tabularx}
\end{center}

Le pipeline B ajoute deux etapes optionnelles : \textbf{alignement} wav2vec2 par langue
(francais natif ; arabe via le modele communautaire \texttt{boualin/wav2vec2-large-xlsr-53-arabic})
et \textbf{diarisation} pyannote (prefixes \texttt{[SPEAKER\_XX]}). La diarisation est
\emph{gated} : sans jeton HuggingFace valide (\texttt{--hf-token}), elle est silencieusement
ignoree et la sortie reste une transcription simple.

% ============================================================
\section{Qualite Darija et adaptateur LoRA}
% ============================================================

Whisper n'a \textbf{pas de code dedie au darija} : le marocain est transcrit sous le code
generique \texttt{ar}, en ecriture arabe. C'est le point faible (WER plus eleve que MSA/
francais). Un comparatif a identifie l'adaptateur LoRA \texttt{anaszil/whisper-large-v3-turbo-darija}
comme la meilleure option.

\begin{center}
\small
\begin{tabular}{lccl}
\toprule
\textbf{Modele} & \textbf{Temps (1h audio)} & \textbf{WER} & \textbf{Note} \\
\midrule
\texttt{anaszil/...-darija} (LoRA) & $\sim$20 s & $\sim$25\,\% & \textbf{le plus rapide et le plus propre} \\
\texttt{MaghrebVoice/...-large-v3} & $\sim$70 s & $\sim$26\,\% & fine-tune complet, 3,4$\times$ plus lent \\
\texttt{large-v3-turbo} (base) & $\sim$15 s & $\sim$28\,\% & sans specialisation darija \\
\bottomrule
\end{tabular}
\end{center}

\textbf{Routage hybride} (\texttt{--darija-lora}) : les chunks detectes \texttt{ar}
passent par la LoRA (PyTorch/PEFT), les chunks \texttt{fr}/\texttt{en} restent sur
CTranslate2 (plus rapide, qualite identique). On ne paie le cout de la LoRA que la ou
elle aide.

% ============================================================
\section{Etat actuel et exemple reel}
% ============================================================

\textbf{Realise :}
\begin{itemize}
  \item Pipeline complet A (decodage $\rightarrow$ VAD $\rightarrow$ chunks $\rightarrow$ detection $\rightarrow$ transcription $\rightarrow$ SRT).
  \item Pipeline B WhisperX (alignement + diarisation optionnels).
  \item Detection de langue par chunk avec allow-list ; routage LoRA Darija.
  \item Tests unitaires du redacteur SRT (round-trip de format).
  \item Notebooks Kaggle (GPU T4) : transcription, WhisperX, comparatif modeles Darija, A/B base vs LoRA.
\end{itemize}

\textbf{Exemple — \texttt{mobachara\_ma3akom\_trimmed}} (talk-show politique, loi de
finances 2026, \og Maroc a deux vitesses \fg{}, hote + 4 invites). La transcription est
\textbf{coherente} : sujets, chiffres et noms correctement captures (\og 380 milliards de
dirhams \fg{}, \og CAN 2026 \fg{}, \og Mondial 2030 \fg{}). La sortie diarisee distingue 5
locuteurs de facon stable. Erreurs darija typiques mais mineures (quelques mots mal
orthographies), sans gener la comprehension.

% ============================================================
\section{Limites et suite}
% ============================================================

\begin{itemize}
  \item \textbf{WER Darija} : point faible structurel ; la LoRA \texttt{anaszil} l'attenue mais ne l'elimine pas.
  \item \textbf{Pas de re-segmentation des sous-titres} : les bornes natives de Whisper sont conservees, d'ou des cues longues (20--24 s) peu confortables pour un affichage humain (acceptables pour le TALN en aval).
  \item \textbf{Diarisation \emph{gated}} : necessite un jeton HuggingFace et l'acceptation des conditions pyannote.
  \item \textbf{Robustesse} : musique, bruit et parole superposee peuvent produire des cues hallucinees ; le \og garble gate \fg{} du benchmark (\texttt{asr\_quality.py}) sert de filet en aval.
  \item \textbf{Travaux futurs} : re-segmentation a budget de caracteres, controle d'un eventuel trou de transcription ($\sim$24 s observe sur l'echantillon), exploration d'un modele Darija dedie.
\end{itemize}

\vspace{4pt}
\rule{\linewidth}{0.4pt}\\[2pt]
{\footnotesize\itshape Sources : \texttt{docs/PIPELINE.md}, \texttt{docs/TRANSCRIPTION.md},
\texttt{docs/WHISPERX\_GUIDE.md}, \texttt{notebooks/} (runners Kaggle et comparatifs).}

% ============================================================
\section{Outils d'ingestion des medias}
% ============================================================

Pour alimenter le pipeline de transcription, deux scripts de telechargement audio ont ete
developpes. Ils constituent la \emph{couche zero} de la chaine : collecter les diffusions
depuis les plateformes sociales (YouTube, Instagram) et les deposer dans une arborescence
normalisee avant que \texttt{transcribe.py} ne les traite.

% ------------------------------------------------------------
\subsection{Telechargeur YouTube incremental — \texttt{fetch\_youtube.py}}
% ------------------------------------------------------------

\textbf{Objectif.} Telecharger en batch l'audio des chaines YouTube de radiodiffuseurs
marocains (2M, Medi1, SNRT, etc.) de maniere \emph{incrementale} : seuls les nouveaux
episodes sont acquis lors des executions suivantes, sans re-telecharger l'existant.

\textbf{Arborescence de sortie.}
\begin{lstlisting}
{compte}/{annee}/{mois}/{titre}.mp3
\end{lstlisting}
Exemple : \texttt{2m\_officiel/2026/07/journal\_20h\_01\_07\_2026.mp3}.
Le nom de fichier est derive du \textbf{titre de la video} (apres nettoyage des
caracteres interdits) plutot que d'un horodatage brut --- ce choix rend l'arborescence
immediatement lisible. L'horodatage \texttt{YYYYMMDDHHMMSS} est calcule en interne
pour determiner les dossiers \texttt{annee/mois} et sert de repli si la video n'a
pas de titre. Cette arborescence permet de retrouver rapidement tout le contenu d'un
mois ou d'un compte sans base de donnees externe.

\textbf{Moteur.} Le telechargement repose sur \textbf{yt-dlp}, invoque avec le format
\texttt{bestaudio/best} et le post-processeur FFmpeg pour la re-encodage en MP3
(192 kbit/s). yt-dlp est prefere a youtube-dl pour sa maintenance active et sa meilleure
gestion des protections anti-bot recentes. La version recente de yt-dlp exige, en plus
d'un interpreteur JavaScript (\texttt{deno} $\geq$ 2.3.0), le paquet Python
\textbf{\texttt{yt-dlp-ejs}} qui fournit le script de resolution des challenges YouTube
en mode hors-ligne, sans necessiter \texttt{--remote-components}. En l'absence de l'un
ou de l'autre, yt-dlp emet un avertissement et certains formats sont indisponibles.

\textbf{Archive de deduplication.} Un fichier d'archive yt-dlp (\texttt{--download-archive})
est maintenu par compte. Il enregistre les identifiants des videos deja acquises ; les
appels subsequents sautent automatiquement ces entrees, ce qui garantit l'idempotence
meme si la liste de la chaine change. Un \emph{filet de securite sur disque} complete
ce mecanisme : si l'archive est absente mais que le fichier audio existe deja sur le
disque, l'archive est \textbf{reconstruite automatiquement} (backfill), evitant ainsi
tout re-telechargement apres la perte accidentelle de l'archive.

\textbf{Progression et journalisation.} Chaque video traitee est prefixee par un
compteur \texttt{[i/N]} (position dans la liste de la chaine) visible en temps reel.
Avant le telechargement effectif, la ligne \texttt{[i/N] [..] downloading <titre>}
s'affiche pour indiquer le debut de l'operation, evitant un silence prolonge sur les
videos volumineuses. Avec \texttt{--log}, chaque ligne est egalement ecrite dans un
fichier journal avec horodatage \texttt{YYYY-MM-DD HH:MM:SS}, compatible avec
\texttt{tail -f}.

\textbf{Options notables.}
\begin{description}[style=nextline, leftmargin=2.2em, labelindent=0.4em]
  \item[\texttt{--scan-limit N}]
    Limite le balayage aux \texttt{N} episodes les plus recents de la chaine (par defaut :
    50). Utile pour les chaines tres actives ou seul le flux recent interesse. Resout
    le blocage observe sur les grandes chaines (ex.~2M TV avec des milliers de videos) :
    sans cette limite, yt-dlp parcourait l'integralite de l'historique avant que le
    filtre \texttt{--max-downloads} ne puisse s'appliquer.
  \item[\texttt{--js-runtimes}]
    Specifie le chemin exact de l'interpreteur JavaScript (\texttt{deno:/chemin/deno},
    \texttt{node}, etc.) a passer a yt-dlp. Necessaire lorsqu'un \texttt{deno} defectueux
    est presente en premier sur le \texttt{PATH} (piege ``deux denos'').
  \item[\texttt{--dry-run}]
    Mode simulation : affiche la liste des episodes qui seraient telecharges sans ecrire
    aucun fichier. Pratique pour valider la configuration avant un vrai batch.
  \item[\texttt{--log}]
    Chemin vers un fichier de journal. Les evenements (nouveau telechargement, saut,
    erreur reseau) y sont horodates ; l'absence de cette option redirige la sortie vers
    \texttt{stdout} uniquement.
  \item[\texttt{--since YYYYMMDD}]
    Ignore les videos publiees avant la date indiquee.
  \item[\texttt{--audio-format}]
    Codec audio cible : \texttt{mp3} (defaut), \texttt{m4a}, \texttt{opus}, etc.
\end{description}

% ------------------------------------------------------------
\subsection{Telechargeur Instagram — \texttt{fetch\_instagram.py}}
% ------------------------------------------------------------

\textbf{Objectif.} Recuperer l'audio des publications et reels Instagram de comptes de
radiodiffuseurs, en suivant la meme arborescence normalisee que pour YouTube.

\textbf{Authentification.} La bibliotheque \textbf{instaloader} est utilisee pour
l'acces a l'API non officielle d'Instagram. La session est chargee depuis un fichier de
cookies ou par identifiants, ce qui evite la re-authentification a chaque execution.

\textbf{Gestion de session et double authentification (2FA).}
Lorsqu'Instagram exige une verification en deux etapes, le script detecte la demande
et invite l'utilisateur a saisir le code TOTP interactivement. La session resultante est
serialisee sur disque (\texttt{--session-file}) pour etre reutilisee sans nouvelle
authentification dans les 90 jours.

\textbf{Extraction audio.} Les publications Instagram sont des videos ; l'audio est extrait
avec \textbf{FFmpeg} en post-traitement, exactement comme pour YouTube, ce qui garantit
un format MP3 uniforme en sortie.

\textbf{Arborescence de sortie.}
\begin{lstlisting}
{compte}/{annee}/{mois}/{titre}.mp3
\end{lstlisting}
Le titre est derive de la \textbf{premiere ligne de la legende} (\emph{caption}) du post
Instagram apres nettoyage des caracteres interdits. En l'absence de legende, le
\textbf{shortcode} du post sert de nom de fichier, garantissant l'unicite.

\textbf{Deduplication.} Un fichier d'archive texte (\texttt{instagram/.download-archive.txt})
memorise les shortcodes deja traites au format \texttt{instagram <shortcode>}, compatible
avec le format yt-dlp. A chaque execution, seuls les nouveaux posts sont telecharges.
Un \emph{filet de securite sur disque} identique a celui du module YouTube reconstruit
automatiquement l'archive si elle est absente.

\textbf{Progression.} Comme pour YouTube, un compteur \texttt{[compte i/N]} prefixe
chaque ligne de sortie, et la ligne \texttt{[..] downloading <titre>} s'affiche avant le
debut effectif du telechargement.

\textbf{Deduplication.} Un registre JSON local (un fichier par compte) memorise les
shortcodes deja traites. A chaque execution, seuls les nouveaux posts sont telecharges.

% ------------------------------------------------------------
\subsection{Telechargeur TikTok -- \texttt{fetch\_tiktok.py}}
% ------------------------------------------------------------

\textbf{Objectif.} Telecharger en batch l'audio des comptes TikTok de radiodiffuseurs
marocains de maniere \emph{incrementale} : seuls les nouveaux clips sont acquis lors
des executions suivantes, sans re-telecharger le contenu deja traite.

\textbf{Arborescence de sortie.}
\begin{lstlisting}
{out}/{compte}/{annee}/{mois}/{titre}.mp3
\end{lstlisting}
Exemple : \texttt{tiktok/2m\_officiel/2026/06/journal-du-soir.mp3}.
Comme pour YouTube, le repertoire racine par defaut est \texttt{./tiktok} ;
il suffit de le pointer vers \texttt{medias/} pour rendre les deux outils
directement interoperables avec \texttt{cli.py}.

\textbf{Moteur yt-dlp (sans moteur JavaScript).} Le script utilise le meme
moteur \textbf{yt-dlp} que \texttt{fetch\_youtube.py}, invoque via l'API Python
\code{YoutubeDL}. A la difference de YouTube, l'extraction TikTok ne necessite
\textbf{aucun interpreteur JavaScript} (Deno, Node, etc.) : yt-dlp resout les
URL TikTok entierement en Python, sans challenge JavaScript cote client. Il n'est
donc pas necessaire d'installer \texttt{yt-dlp-ejs} ni un runtime externe pour
ce script.

\textbf{Listage en une seule phase.} Contrairement au listage YouTube en deux
phases (liste plate puis metadonnees par video), l'API TikTok retourne
\textbf{deja le timestamp, le titre et l'identifiant de l'auteur} dans chaque
entree de la liste plate (\texttt{extract\_flat="in\_playlist"}). Un seul appel
yt-dlp suffit donc pour obtenir toutes les metadonnees necessaires au filtrage
et a la construction du chemin de sortie, sans appel complementaire par video.
Cela rend le listage incrementiel particulierement rapide sur les grands comptes.

\textbf{Deduplication.} Le meme mecanisme en deux couches que les autres outils
est applique :
\begin{itemize}
  \item \textbf{Archive de telechargement} : un fichier texte
    \texttt{.download-archive.txt} a la racine de sortie enregistre une entree
    \texttt{tiktok <id>} par clip telecharge ; les executions suivantes sautent
    immediatement les entrees deja presentes.
  \item \textbf{Filet de securite sur disque} : si le fichier \texttt{.mp3} cible
    existe deja sans entree d'archive correspondante (ex.~apres une perte de
    l'archive), l'archive est reconstruite (\emph{backfill}) et le clip est
    ignore, evitant tout re-telechargement inutile.
\end{itemize}

\textbf{Comptes prives (\texttt{--cookies-file}).} Pour acceder a des comptes
TikTok prives, l'option \texttt{--cookies-file} permet de fournir un fichier de
cookies (au format Netscape, exportable depuis le navigateur). Ce fichier est
transmis directement a yt-dlp pour authentifier les requetes, sans stocker de
mot de passe en clair.

\textbf{Options notables.}
\begin{description}[style=nextline, leftmargin=2.2em, labelindent=0.4em]
  \item[\texttt{--account}]
    Identifiant (handle) du compte TikTok a telecharger. Repetable pour traiter
    plusieurs comptes en une seule invocation
    (ex.~\texttt{--account 2m --account medi1}).
  \item[\texttt{--scan-limit N}]
    Limite le balayage aux \texttt{N} clips les plus recents du compte (defaut :
    50). Passer \texttt{0} desactive la limite et parcourt tout l'historique.
  \item[\texttt{--since YYYYMMDD}]
    Ignore les clips publies avant la date indiquee.
  \item[\texttt{--max-downloads N}]
    Plafonne le nombre de nouveaux telechargements pour cette execution.
  \item[\texttt{--dry-run}]
    Mode simulation : affiche la liste des clips qui seraient telecharges sans
    ecrire aucun fichier.
  \item[\texttt{--cookies-file}]
    Chemin vers un fichier de cookies Netscape ; necessaire pour les comptes
    prives.
  \item[\texttt{--log}]
    Chemin vers un fichier journal ; les evenements y sont horodates
    (\texttt{YYYY-MM-DD HH:MM:SS}).
  \item[\texttt{--audio-format}]
    Codec audio cible : \texttt{mp3} (defaut), \texttt{m4a}, \texttt{opus}, etc.
\end{description}

% ------------------------------------------------------------
\subsection{Module de fonctions partage — \texttt{\_media\_common.py}}
% ------------------------------------------------------------

Les deux scripts s'appuient sur un module interne \texttt{\_media\_common.py} qui
centralise la logique reutilisable :

\begin{itemize}
  \item Construction et validation de l'arborescence de sortie
    (\texttt{\{compte\}/\{annee\}/\{mois\}/}).
  \item Normalisation des noms de fichiers (suppression des caracteres interdits,
    troncature des titres longs, encodage UTF-8 coherent sous Linux et Windows).
  \item Fonctions d'horodatage et de journalisation (format partage entre les deux scripts).
  \item Gestion des tentatives avec repli exponentiel (\emph{exponential back-off}) pour
    les erreurs reseau transitoires.
  \item Utilitaires FFmpeg : detection du binaire, construction de la commande de
    conversion, verification du fichier de sortie.
\end{itemize}

Ce module n'a pas d'entree CLI propre ; il est importe par \texttt{fetch\_youtube.py}
et \texttt{fetch\_instagram.py}.

% ------------------------------------------------------------
\subsection{Suite de tests}
% ------------------------------------------------------------

Les quatre modules (\texttt{fetch\_youtube.py}, \texttt{fetch\_instagram.py},
\texttt{fetch\_tiktok.py} et \texttt{\_media\_common.py}) sont couverts par une suite de
\textbf{216 tests} executes avec \textbf{pytest}. Les 38 nouveaux tests TikTok
(\texttt{test\_fetch\_tiktok.py}) utilisent un bouchon \texttt{FakeTikTokYDL} qui
intercepte les appels yt-dlp sans connexion reseau. Les cas couverts incluent la
construction de l'arborescence, la deduplication, la gestion des erreurs, et le mode
\texttt{--dry-run}. L'ensemble de la suite reste reproductible sans connexion reseau
ni media reels.

\vspace{4pt}
\rule{\linewidth}{0.4pt}\\[2pt]
{\footnotesize\itshape Sources : \texttt{ingestion/fetch\_youtube.py},
\texttt{ingestion/fetch\_instagram.py}, \texttt{ingestion/\_media\_common.py},
\texttt{tests/} (suite pytest, 178 tests).}

\end{document}
